% Generated by Listas/make_statements.py from pratica-lecture-07.md, item 21. Do not edit.
\textbf{Source.} Tadelis (2013), \emph{Game Theory: An Introduction}, Exercise 12.7 ``Trading Places'', p.~268 {[}PDF p.~287{]}. Reproduced verbatim. \textbf{Answer in the source:} no (none in the book; Princeton UP's website has about half, not checked).



\begin{quote}

\textbf{Trading Places:} Two players, 1 and 2, each own a house. Each player \(i\) values his own house at \(v_i\). The value of player \(i\)'s house to the other player, i.e., to player \(j \ne i\), is \(\frac{3}{2}v_i\). Each player \(i\) knows the value \(v_i\) of his own house to himself, but not the value of the other player's house. The values \(v_i\) are drawn independently from the interval \([0, 1]\) with uniform distribution.



\begin{enumerate}

\def\labelenumi{\alph{enumi}.}

\item

  Suppose players announce simultaneously whether they want to exchange their houses. If both players agree to an exchange, the exchange takes place. Otherwise no exchange takes place. Find a Bayesian Nash equilibrium of this game in pure strategies in which each player \(i\) accepts an exchange if and only if the value \(v_i\) does not exceed some threshold \(\theta_i\).

\item

  How would your answer to (a) change if player \(j\)'s valuation of player \(i\)'s house were \(\frac{5}{2}v_i\)?

\item

  Try to explain why any Bayesian Nash equilibrium of the game described in (a) must involve threshold strategies of the type postulated in (a).

\end{enumerate}

\end{quote}

